\section{The sphere-null counterexample}

We write $S^n_{(2)}$ for the nilpotent $2$-localization of $S^n$, and set
\[
 \Nul=\{X\in\cS_*:\Map_*(X,S^n_{(2)})\simeq *\text{ for all }n\geq2\}.
\]
The mapping space in this definition includes all its components. Vanishing only of the positive homotopy groups of its constant component would be a weaker condition.

\begin{claim}{A.1}{Proved}{Desuspension, including the bottom dimension}\label{A1}
For every pointed CW complex $X$, $\Sigma X\in\Nul$ implies $X\in\Nul$.
\end{claim}
\begin{proof}
The precise classical input is the $2$-local EHP homotopy fibre sequence
\[
 S^n_{(2)}\xrightarrow{E}\Omega S^{n+1}_{(2)}
       \xrightarrow{H}\Omega S^{2n+1}_{(2)},\qquad n\geq1.
\]
Here $E$ is adjoint to suspension and $H$ is the James--Hopf map. We opened its statement in \citet[Theorem 2.2]{Boyde2023}, whose convention immediately before the theorem is that all spaces are $2$-local. The same statement, including the even-dimensional case, appears in \citet[Theorem 1(ii), with $p=2$]{RavenelEHP2005}. We use an actual homotopy fibre sequence in all dimensions, not merely the metastable exact sequence of homotopy groups.

Apply the derived pointed mapping-space functor $\Map_*(X,-)$. It preserves homotopy fibres. Its values on the second and third terms are
\[
 \Map_*(\Sigma X,S^{n+1}_{(2)}),\qquad
 \Map_*(\Sigma X,S^{2n+1}_{(2)}),
\]
both weakly contractible by hypothesis for every $n\geq2$. Their homotopy fibre is consequently weakly contractible. For $n=2$ the sequence is specifically
$S^2_{(2)}\to\Omega S^3_{(2)}\to\Omega S^5_{(2)}$; no dimensional exception or limiting argument is involved. This proves the assertion simultaneously on components and on every homotopy group.
\end{proof}

\begin{claim}{A.2}{Proved}{The closure properties actually needed}\label{A2}
The class $\Nul$ is invariant under pointed weak equivalence, contains the point, and is closed under retracts, all pointed homotopy colimits, and all three two-out-of-three implications in a pointed homotopy cofiber sequence. Moreover $X\in\Nul$ if and only if $\Sigma X\in\Nul$. Every member is connected.
\end{claim}
\begin{proof}
For a pointed diagram $X_i$, mapping out of its homotopy colimit gives the homotopy limit of $\Map_*(X_i,S^n_{(2)})$. A diagram of weakly contractible spaces has weakly contractible homotopy limit. This proves the colimit assertion; invariance and retracts are immediate from the same mapping-space definition. Suspension is a pointed homotopy colimit, and its converse is A.1.

For a cofiber sequence $X\to Y\to C$, apply $\Map_*(-,S^n_{(2)})$. If $X,Y\in\Nul$, its fibre sequence shows $C\in\Nul$; if $X,C\in\Nul$, it shows $Y\in\Nul$. If $Y,C\in\Nul$, apply the first implication to the next cofiber sequence $Y\to C\to\Sigma X$ and then A.1. This last use of desuspension is indispensable.

If $X$ has a component not containing its basepoint, choose a point in that component. Evaluation there has a section $S^2_{(2)}\to\Map_*(X,S^2_{(2)})$, obtained by making a map constant on that component and equal to the basepoint on every other component. Components of a CW complex are open and closed, so these maps are continuous. The noncontractible space $S^2_{(2)}$ is thus a retract of this mapping space, contradicting $X\in\Nul$.
\end{proof}

\begin{claim}{A.3}{Proved}{The cofiber construction gives a reversible fundamental localizer}\label{A3}
Let $\mathcal A$ be a weak-equivalence-invariant class of pointed spaces containing the point, closed under pointed homotopy colimits, retracts, and two out of three in homotopy cofiber sequences. Define
\[
 W_{\mathcal A}=\{u:A\to B\text{ in }\Cat:C(Bu)\in\mathcal A\},
 \qquad C(f)=Y\amalg_X^h *.
\]
Then $W_{\mathcal A}$ is a fundamental localizer in the relative-Theorem-A convention and is reversible. Its aspherical nonempty spaces are exactly the underlying pointed spaces in $\mathcal A$. In particular this applies to $\mathcal A=\Nul$, giving $W_S$.
\end{claim}
\begin{proof}
The cofiber of an identity is contractible. The pointed pushout-pasting law gives a cofiber sequence
$C(f)\to C(gf)\to C(g)$ for composable unpointed maps. One can see it by applying the pointed cofiber construction to $f_+$ and $g_+$ and collapsing the first cone inside the iterated cone. Thus two out of three in $\mathcal A$ gives two out of three for $W_{\mathcal A}$.

Functoriality of homotopy cofiber sends an arrow retract to a pointed homotopy retract, so $W_{\mathcal A}$ is closed under arrow retracts. This directly checks the retraction rule in the proposed note. More explicitly, if $i:X\to Y$, $r:Y\to X$, $ri=\id_X$ and $ir\in W_{\mathcal A}$, then the arrow $r:Y\to X$ is a retract of $ir:Y\to Y$: on sources both maps are identities and on targets they are $i$ and $r$. Hence $r\in W_{\mathcal A}$, and $i\in W_{\mathcal A}$ follows from $ri=\id_X$. No asserted splitting of a suspended cofiber is needed.

A terminal object contracts the nerve of its category, checking the terminal-object axiom. For relative Theorem A take a triangle $A\xrightarrow{u}B\to D$. The diagrams $d\mapsto A/d$ and $d\mapsto B/d$ are covariant, and there are natural weak equivalences
\[
 \hocolim_{d\in D}B(A/d)\simeq BA,
 \qquad \hocolim_{d\in D}B(B/d)\simeq BB.
\]
For completeness the projection from the Grothendieck construction of $d\mapsto A/d$ to $A$ has the left adjoint $a\mapsto(v(a),a,\id_{v(a)})$. The nerve homotopy-colimit comparison therefore gives the first displayed equivalence, and likewise the second, naturally in $u$. The homotopy cofiber of the induced map is the \emph{pointed} homotopy colimit of the cofibers $C(B(u/d))$. To check the basepoints, apply $(-)_+$ to both diagrams and then take their pointed cofibers: both operations preserve colimits. Equivalently the section of cone vertices in the unpointed construction is collapsed. If every $u/d$ belongs to $W_{\mathcal A}$, closure of $\mathcal A$ under pointed homotopy colimits now proves $u\in W_{\mathcal A}$.

The cofiber of $X\to *$ is its unreduced suspension; for nonempty pointed $X$ it is pointed-homotopy-equivalent to its reduced suspension. The equivalence $X\in\mathcal A\Longleftrightarrow\Sigma X\in\mathcal A$ follows by applying cofiber two out of three to $X\to *\to\Sigma X$. Hence $u\in W_{\mathcal A}$ if and only if $C(Bu)$ is $W_{\mathcal A}$-aspherical. For a homotopy-pushout square the cofibers of its two parallel edges are naturally equivalent by pushout pasting. The preceding criterion therefore gives the reverse sieve-cover implication.

This also checks the precise part of POLES B.2 and B.5 used here. For a cover by two full sieves the nerve is literally the union of their nerves over the nerve of the intersection: every string of arrows lies in the piece containing its last vertex. The resulting pushout is a homotopy pushout because both inclusions are simplicial monomorphisms. In the converse direction, POLES's category $K(u)$ is the two-layer mapping collage with one cone point on the source layer; its two full sieves are the mapping collage and the source cone. The mapping collage retracts onto the target by an adjunction. Applying the reverse sieve implication to that square gives the cofiber criterion, with the actual functor $u$ recovered by this retraction.

If the source of $f$ is empty, then $C(f)=Y_+$. If both source and target are empty, this is the one-point space. The construction and relative-Theorem-A argument include these cases. For the class $\Nul$, a nonempty $Y_+$ is disconnected and so does not belong to $\Nul$ by A.2; in particular the empty category has not inadvertently become aspherical.
\end{proof}

\begin{claim}{A.4}{Proved}{The unstable separating witnesses}\label{A4}
We have $\RP^\infty\in\Nul$ and $\RP^{2k}\notin\Nul$ for every $k\geq1$. Consequently $\RP^\infty\to *$ belongs to $W_S$, whereas $\RP^2\to *$ does not, and $W_S$ is nontrivial.
\end{claim}
\begin{proof}
The exact input from the attached Miller paper is Theorem C on page 40: if a nilpotent space $Z$ has bounded mod-$p$ cohomology, then $\Map_*(B\mathbb Z/p,Z)$ is weakly contractible \citep[Theorem C]{Miller1984}. Each $S^n_{(2)}$, $n\geq2$, is simply connected and has mod-$2$ cohomology only in degrees $0,n$. The theorem applies directly; an arithmetic-square argument is unnecessary. Since $B\mathbb Z/2\simeq\RP^\infty$, this proves the first assertion.

For a CW complex of dimension $d$ and a $(d-1)$-connected target $Z$, cellular obstruction theory identifies based homotopy classes into $Z$ with $H^d(-;\pi_d Z)$. Indeed maps on the $(d-1)$-skeleton are uniquely null up to the required lower homotopies; choices on $d$-cells are a cellular cocycle, and homotopies change it by a coboundary. With $d=2k$ this gives
\[
 [\RP^{2k},S^{2k}_{(2)}]_*\cong
 H^{2k}(\RP^{2k};\mathbb Z_{(2)})\cong\mathbb Z/2.
\]
The last calculation uses the cellular cochain differential $2$ into the top even cell. It is nonzero. A.2 and A.3 translate these two membership tests into the claimed functors. Independently the identity of $S^2_{(2)}$ shows that $S^2_{(2)}$ is not aspherical.
\end{proof}

\begin{claim}{A.5}{Proved}{The stable subcategory generated by the infinite projective space}\label{A5}
In the stable homotopy category,
\[
 \loc(\Sigma^\infty\RP^\infty)=\loc(\mathbb S/2).
\]
Every reversible fundamental localizer containing $\RP^\infty\to *$ therefore contains $\RP^2\to *$ after applying $RT$. In particular $W_S\ne RT(W_S)$.
\end{claim}
\begin{proof}
Put $V=\Sigma^\infty\RP^\infty$ and $\mathcal L=\loc(V)$. Every localizing subcategory of spectra is a tensor ideal: for fixed $V'\in\mathcal L$, the spectra $E$ with $V'\wedge E\in\mathcal L$ form a localizing subcategory containing the sphere, and sphere cells generate spectra. This confirms the step of POLES F.13 used below.

The Snaith splitting we use is
\[
 \Sigma^\infty QX\simeq
 \bigvee_{j\geq1}(\Sigma^\infty X)^{\wedge j}_{h\Sigma_j}
 \quad (X\text{ connected and well pointed}).
\]
We opened \citet[Theorem 1.1 and the filtration of Section 2]{Snaith1974}. Its $n=\infty$ configuration space in degree $j$ is contractible with free $\Sigma_j$ action; its $j$th filtration quotient is thus $E\Sigma_{j+}\wedge_{\Sigma_j}X^{\wedge j}$. This identifies the theorem's summands with those in the displayed formula. The tensor-ideal and homotopy-colimit properties give $\Sigma^\infty Q\RP^\infty\in\mathcal L$.

The Kahn--Priddy input is a \emph{section of spaces}. Its opened formulations are \citet[Theorem 3.1]{KahnPriddy1972}, the unnumbered main theorem on page 103 of \citet{KahnPriddy1978}, and \citet[Theorem 7.5]{Snaith2009}: if $f:\RP^\infty\to Q_0S^0$ induces an isomorphism on $\pi_1$, its infinite-loop extension $\widetilde f:Q\RP^\infty\to Q_0S^0$ admits a map in the opposite direction whose composite with $\widetilde f$ is a $2$-local equivalence. Localizing and composing with a homotopy inverse makes $(Q_0S^0)_{(2)}$ a retract of $(Q\RP^\infty)_{(2)}$ as a pointed space. We neither assert nor need an infinite-loop section. Localization of these connected simple spaces induces localization of their reduced suspension spectra, so
\[
 \Sigma^\infty(Q_0S^0)_{(2)}\in\mathcal L.
\]

We spell out the passage from this space to a spectrum. If a connected grouplike $H$-space $Z$ has $\Sigma^\infty Z\in\mathcal L$, the same is true of its bar construction $BZ$: the reduced suspension spectrum of each $Z^r$ is a finite wedge of smash products of $\Sigma^\infty Z$, and the bar realization is a homotopy colimit. This remains true at every iterated bar stage for an infinite-loop space. For $E=\tau_{\geq1}\mathbb S_{(2)}$, these stages are $\Omega^\infty\Sigma^rE$. There is a canonical equivalence
\[
 E\simeq\hocolim_{r\geq0}\Sigma^{-r}\Sigma^\infty
                       \Omega^\infty\Sigma^rE.
\]
For a fixed homotopy group the comparison becomes an isomorphism once $r$ is large, by Freudenthal applied to the increasingly connected spaces $\Omega^\infty\Sigma^rE$; the maps are the adjoints of the spectrum structure maps. Thus $E\in\mathcal L$.

The same bar argument starting with $B\mathbb Z/2$ gives all $K(\mathbb Z/2,r)$ and then, by the displayed stabilization construction, $H\mathbb Z/2\in\mathcal L$. Smash the canonical triangle
$\tau_{\geq1}\mathbb S_{(2)}\to\mathbb S_{(2)}\to H\mathbb Z_{(2)}$
with $\mathbb S/2$. Its first and third terms belong to $\mathcal L$; hence so does its middle term $\mathbb S/2$.

For the reverse inclusion, $V[1/2]=0$. Indeed its integral homology is $2$-primary torsion, and a connective spectrum with zero integral homology is zero: its first nonzero homotopy group, if any, would be detected by the stable Hurewicz theorem. The kernel of $(-)[1/2]$ is $\loc(\mathbb S/2)$. One explicit verification is that $\mathbb S/2^r$ is built by extensions of $\mathbb S/2$, while the fibre of $E\to E[1/2]$ is the telescope of the spectra $\Sigma^{-1}E/2^r$. Tensor-ideal closure then gives every object in that kernel. This proves the equality.

Finally, if $W$ is reversible and contains $\RP^\infty\to *$, then $T(W)$ contains $V$, since the cofiber of this map is a suspension of $\RP^\infty$. It therefore contains $\mathbb S/2$ and the reduced suspension spectrum $\Sigma^\infty\RP^2\simeq\Sigma(\mathbb S/2)$. The definition of $R$ implies $\RP^2\to *\in RT(W)$. This confirms the specific Galois step of POLES F.4--F.6: $T(W)\subseteq\mathcal L$ is equivalent, by the generators in its definition, to $W\subseteq R(\mathcal L)$; the spectrum assigned to a functor is exactly the reduced suspension spectrum of its unpointed homotopy cofiber. A.4 supplies the separating nonmembership for $W_S$.
\end{proof}

\begin{claim}{A.6}{Proved}{The odd-prime desuspension repair}
For an odd prime $p$ put $J_{p-1}(S^{2n})=\widehat S^{2n}$ and let $\mathcal Z_p$ consist of all $S^{2n-1}_{(p)}$ and all $\widehat S^{2n}_{(p)}$, $n\geq1$. Its null class $\Nul_p$ is closed under desuspension. The resulting reversible localizer contains $B\mathbb Z/p\to *$ and excludes $M(\mathbb Z/p,1)\to *$.
\end{claim}
\begin{proof}
The two Toda fibre sequences, with all spaces $p$-localized, are
\[
 \widehat S^{2n}\longrightarrow\Omega S^{2n+1}
       \longrightarrow\Omega S^{2np+1},\qquad
 S^{2n-1}\longrightarrow\Omega\widehat S^{2n}
       \longrightarrow\Omega S^{2np-1}.
\]
We opened their explicit statement in \citet[p.~80, item (ii)]{BodigheimerHenn1983}; the first also appears as Theorem 1(ii) of \citet{RavenelEHP2005}. Apply $\Map_*(X,-)$ exactly as in A.1. If $\Sigma X$ is null for the entire displayed family, both loop-space terms of each sequence have contractible mapping space, proving nullity of each fibre. The $n=1$ case includes $S^1_{(p)}$ and $J_{p-1}(S^2)_{(p)}$ and therefore supplies the bottom step explicitly.

Each target is nilpotent and has bounded mod-$p$ cohomology: the James stage has finitely many cells before localization, and an odd sphere, including the circle, has bounded cohomology. Miller's Theorem C therefore makes $B\mathbb Z/p$ null for all of them. The target $J_{p-1}(S^2)_{(p)}$ is simply connected with $\pi_2\cong\mathbb Z_{(p)}$, so obstruction theory on the two-dimensional Moore space gives
\[
 [M(\mathbb Z/p,1),J_{p-1}(S^2)_{(p)}]_*
      \cong H^2(M(\mathbb Z/p,1);\mathbb Z_{(p)})
      \cong\mathbb Z/p.
\]
It is not null. The rest of A.2--A.3 applies unchanged. This is the precise modification required by the Toda sequences: at an odd prime their even-dimensional fibre is a truncated James construction, not the even sphere. We do not replace that fibre by a sphere without an additional argument.
\end{proof}

\begin{claim}{A.7}{Proved}{The stable comparison at every prime}
For every prime $p$,
\[
 \loc(\Sigma^\infty B\mathbb Z/p)=\loc(\mathbb S/p).
\]
Consequently the odd-prime localizer of A.6 is not $RT$-fixed, and its bounded-subcomplex version of B.2 gives an accessible such example.
\end{claim}
\begin{proof}
Put $\mathcal L_p=\loc(\Sigma^\infty B\mathbb Z/p)$. We justify the only new input to A.5. The exact all-prime splitting statement in \citet[Theorem F]{CarusoCohenMayTaylor1984} supplies maps of pointed spaces
\[
 Q_0S^0\xrightarrow{j'_p}QB\Sigma_p
        \xrightarrow{s'_p}Q_0S^0
\]
whose composite is a $p$-local equivalence. The section $j'_p$ is not asserted to be an infinite-loop map. The original paper \citep[unnumbered main theorem, p.~103]{KahnPriddy1978} states the same splitting and explicitly permits $B\mathbb Z/p$ in place of $B\Sigma_p$.

Here is a derivation using the displayed form, which also checks the change of group. The covering $B\mathbb Z/p\to B\Sigma_p$ has degree $(p-1)!$. Its stable transfer, followed by the covering projection, induces multiplication by this degree on ordinary homology; the compatibility of the stable and ordinary transfers is \citet[Proposition 2.1 and Remark 2.3]{KahnPriddy1972}. The maps on reduced suspension spectra are obtained by removing their sphere summands. After $p$-localization this composite is an integral-homology equivalence of connective spectra and hence a stable equivalence, by the stable Hurewicz argument of A.5. Thus $(\Sigma^\infty B\Sigma_p)_{(p)}$ is a retract of $(\Sigma^\infty B\mathbb Z/p)_{(p)}$ and belongs to $\mathcal L_p$. We have used a homology equivalence of the composite, not the stronger and unnecessary assertion that the composite is multiplication by the degree as a stable map.

The Snaith splitting of A.5 now puts $(\Sigma^\infty QB\Sigma_p)_{(p)}$ in $\mathcal L_p$. The displayed section puts $\Sigma^\infty(Q_0S^0)_{(p)}$ there. Iterated bars and stabilization give $\tau_{\geq1}\mathbb S_{(p)}\in\mathcal L_p$. Iterated bars of $B\mathbb Z/p$ give $H\mathbb Z/p\in\mathcal L_p$. Smashing the connective-cover triangle of $\mathbb S_{(p)}$ with $\mathbb S/p$ then puts $\mathbb S/p$ in $\mathcal L_p$. Conversely $(\Sigma^\infty B\mathbb Z/p)[1/p]=0$, so the telescope argument of A.5 gives the reverse inclusion. Finally $\Sigma^\infty M(\mathbb Z/p,1)\simeq\Sigma(\mathbb S/p)$, and A.6 excludes this Moore space from the unstable aspherical class. The same countable witnesses belong, or fail to belong, to its $\aleph_1$ bounded-subcomplex version.
\end{proof}
